% ============================================================
% paper/appendices/appx_minimal_oc_gateway.tex
% Appendix C: Minimal Ontological Gateway (OC / K7)
% Purpose: minimal non-exhaustive definitions sufficient to read K_EA.
% ============================================================

\section*{Appendix C: Minimal Ontological Gateway (OC / K7)}

\subsection*{Purpose of This Appendix}

This appendix provides a minimal gateway into the Ontology of Continua
(OC) context sufficient to interpret the usage of $K_{\mathrm{EA}}$ and
its parent coupling to $K7$ as presented in this whitepaper.

It is explicitly \textbf{non-exhaustive}. It does not restate OC Core,
does not introduce new primitives, and does not modify any executable
content of \texttt{ETS.K\_EA.v1.1}.

\subsection*{Minimal OC Vocabulary (Non-Exhaustive)}

The whitepaper uses a compact OC-style vocabulary. The intended reading
is purely structural.

\begin{definition}[Continuum]
A \emph{continuum} $K$ denotes a closed formal object whose admissible
states are represented by a set $\Omega(K)$, together with explicit
boundary conditions that constrain the admissible region.
\end{definition}

\begin{definition}[Admissible state space]
$\Omega(K)$ denotes the set of admissible states of $K$ under the
constraints and thresholds declared for that continuum. If
$\Omega(K)=\varnothing$, the continuum has no admissible states under the
declared constraints.
\end{definition}

\begin{definition}[Boundary]
$\partial \Omega(K)$ denotes the boundary of admissibility: states at
which one or more admissibility conditions are met at equality, or where
stability is lost. In this document, boundary language is used as a
structural delimiter, not as a physical claim.
\end{definition}

\begin{definition}[Axes and degrees of freedom]
$A(K)$ denotes the set of axes (dimensions of distinguishability) used to
parameterize or describe admissible states. The term \emph{degree of
freedom} is used descriptively: it refers to the presence of an axis in
the continuum representation, not to any specific numeric count imposed
across enterprises.
\end{definition}

\begin{definition}[Thresholds]
$\Theta(K)$ denotes the set (or vector) of threshold conditions that
separate admissible from non-admissible regimes. Thresholds in this
whitepaper are structural boundaries of the formal object and are not
presented as empirically universal constants.
\end{definition}

\begin{definition}[Flows and cycles]
$J(K)$ denotes admissible flows (channels of state change) and $C(K)$
denotes closure/cycle structures required by the continuum definition.
These symbols are used as structural placeholders for closure logic
defined by the ETS.
\end{definition}

\begin{definition}[Continuity measure]
$k(K)$ denotes a continuity measure used to summarize whether the
continuum retains coherent admissibility under the declared constraints.
For the enterprise continuum, this is written as $k_{\mathrm{EA}}$.
\end{definition}

\subsection*{Minimal Interpretation of \texorpdfstring{$K7$}{K7}}

In this whitepaper, $K7$ is used as a \emph{parent context layer} for
$K_{\mathrm{EA}}$ to express coupling to a broader social/institutional
environment. The usage is intentionally minimal:

\begin{itemize}
  \item $K_{\mathrm{EA}}$ is treated as a continuum whose admissibility
  may depend on external constraints not internal to the enterprise.
  \item Those external constraints are represented abstractly as a
  parent context (denoted $K7$), without requiring a full OC exposition.
  \item The coupling is descriptive: it states that compatibility and
  observability conditions may be conditioned by a parent layer.
\end{itemize}

This appendix does not define a complete theory of $K7$.
It provides only the minimal statement needed to interpret references to
``Parent Coupling (K7)'' in the main text.

\subsection*{Non-Claims}

This gateway makes no claims regarding:
(i) completeness of OC,
(ii) empirical correctness of any layer interpretation,
(iii) uniqueness of the chosen symbols, or
(iv) usefulness of layering for practice.

Its function is to prevent reader opacity by supplying minimal
definitions used by the text.

%%%End of Document%%%
